\section{不变方向与特征值：寻找矩阵最自然的坐标轴}
\label{sec:03-Linear-Algebra/05-Eigenvalues-and-Eigenvectors/01}

随便拿一个矩阵作用在向量上，长度和方向通常都跟着变了。连续作用 \(A,A^2,A^3,\ldots\)，坐标分量反复混合，很快就面目全非——直接算下去既笨又没有结构感。

但如果你能在这个变换里找到一些“无论怎么作用都不转向”的特殊方向，整件事就一下子清晰了。沿这些方向，\(A\) 只做最简单的操作——拉伸、压缩或翻转，不搅动任何其他分量。

这就是特征向量的出发点。

\subsection{从定义到方程}

非零向量 \(v\) 如果满足

\[
Av = \lambda v,
\]

那么 \(A\) 对它的作用仅仅是缩放 \(\lambda\) 倍。\(v\) 叫特征向量，\(\lambda\) 叫对应的特征值。

要求 \(v\neq 0\) 是关键的。零向量对任何 \(\lambda\) 都满足等式 \(A\cdot0=\lambda\cdot0\)，但它没有任何方向信息——你没办法拿一个零向量去指认“这里是我不转的方向”。

把定义移项：

\[
(A-\lambda I)v = 0.
\]

我们需要这个齐次系统有非零解。这就意味着 \(A-\lambda I\) 必须不可逆——否则只有零解。用行列式来判断可逆性：

\[
\det(A-\lambda I) = 0.
\]

左边是关于 \(\lambda\) 的 \(n\) 次多项式，叫特征多项式。它的根就是特征值。对每一个求出来的根，再解齐次系统

\[
\mathcal N(A-\lambda I)
\]

得到对应的特征空间。

这条推导也解释了为什么计算步骤总是“先求特征值，再求特征向量”：只有先在 \(A-\lambda I\) 奇异的条件下，非零方向才有存在的余地。

\subsection{一个能直接看见的例子}

取

\[
A = \begin{pmatrix}2&1\\1&2\end{pmatrix}.
\]

动手之前先凭眼睛找找看。尝试 \((1,1)^{\trans}\)：

\[
A\begin{pmatrix}1\\1\end{pmatrix}
= \begin{pmatrix}2+1\\1+2\end{pmatrix}
= \begin{pmatrix}3\\3\end{pmatrix}
= 3\begin{pmatrix}1\\1\end{pmatrix}.
\]

向量方向没动，只是放大了三倍。再试 \((1,-1)^{\trans}\)：

\[
A\begin{pmatrix}1\\-1\end{pmatrix}
= \begin{pmatrix}2-1\\1-2\end{pmatrix}
= \begin{pmatrix}1\\-1\end{pmatrix}.
\]

这个方向连缩放都没有——进去什么样，出来还是什么样。

这个矩阵在标准坐标中的两个分量互相耦合：原来 \((x,y)\) 变成 \((2x+y,x+2y)\)，你一眼看不出规律。但如果沿 \((1,1)\) 和 \((1,-1)\) 这两条斜轴去看，变换简单到只是独立的缩放。矩阵本身没有变——换的是你看它的角度。

从特征多项式也能得到同样的结果：

\[
\begin{aligned}
\det(A-\lambda I)
&= \det\begin{pmatrix}2-\lambda&1\\1&2-\lambda\end{pmatrix} \\
&= (2-\lambda)^2 - 1 \\
&= (\lambda-3)(\lambda-1).
\end{aligned}
\]

根为 \(\lambda=3\) 和 \(\lambda=1\)，正是我们肉眼找到的那两个。

\subsection{缩放因子的不同含义}

若 \(Av=\lambda v\)，\(\lambda\) 可正可负可为零：

\begin{itemize}
\tightlist
\item
  \(\lambda>1\)：沿该方向放大——作用一次拉长，作用多次指数增长；
\item
  \(0<\lambda<1\)：沿该方向缩小但不翻转——越来越短但不改变指向；
\item
  \(\lambda<0\)：缩放同时翻转方向——每次作用都会掉头；
\item
  \(\lambda=0\)：该方向被压到零——所有位于该特征空间的向量都落入零空间，因此 \(v\in\mathcal N(A)\)；
\item
  \(|\lambda|=1\)：沿该方向长度保持。但仅凭这一条不足以判断长期行为——还需要看特征方向是否足够。第三节会给出 \(|\lambda|=1\) 但系统仍在增长的例子。
\end{itemize}

零是特征值当且仅当 \(A\) 不可逆。理由很直接：

\[
Av=0,\;v\neq0
\;\Longleftrightarrow\;
\mathcal N(A)\neq\{0\}
\;\Longleftrightarrow\;
\det A=0.
\]

特征值 \(\lambda=0\) 的存在意味着存在非零向量被整个拍到原点——这是不可逆的几何翻译。

用一个秩 \(1\) 的矩阵快速体会一下：

\[
A=\begin{pmatrix}1&2\\2&4\end{pmatrix}.
\]

第二列是第一列的两倍，列线性相关，矩阵不可逆。特征多项式：

\[
\det(A-\lambda I)=\det\begin{pmatrix}1-\lambda&2\\2&4-\lambda\end{pmatrix}
= (1-\lambda)(4-\lambda)-4 = \lambda^2-5\lambda = \lambda(\lambda-5).
\]

根为 \(\lambda_1=0,\;\lambda_2=5\)。\(\lambda_1=0\) 的特征向量满足 \((A-0I)v=0\)，也就是 \(v_1+2v_2=0\)，所以 \((2,-1)^{\trans}\) 及其倍数就是零空间。这些向量被 \(A\) 整个拍扁到原点。\(\lambda_2=5\) 的方向是 \((1,2)^{\trans}\)——放大五倍。

迹是 \(1+4=5\)，恰好是 \(\lambda_1+\lambda_2=0+5\)。行列式是 \(1\cdot4-2\cdot2=0\)，恰好是 \(\lambda_1\cdot\lambda_2=0\cdot5\)。矩阵不可逆的等价描述又多了一条：至少有一个特征值为零。

\subsection{特征空间为什么是子空间}

固定一个特征值 \(\lambda\)，所有对应的特征向量再补上零向量，构成

\[
E_\lambda = \mathcal N(A-\lambda I).
\]

这个集合是子空间。若 \(u,v\in E_\lambda\)，则

\[
A(\alpha u+\beta v) = \lambda(\alpha u+\beta v),
\]

线性组合仍在同一个特征空间里。

同一特征值可能对应多个线性无关的特征方向。\(E_\lambda\) 的维数叫几何重数——也就是真正能找到几个独立的不转方向。特征值作为特征多项式根的出现次数叫代数重数——来自代数方程的重根计数。

两者之间有一个恒成立的不等式：

\[
1 \leq \dim E_\lambda \leq \text{代数重数}.
\]

下界是显然的：每个特征值按定义至少带来一个特征向量。上界可以这样看：取 \(E_\lambda\) 的一组基扩充为全空间基，在这组基下矩阵的前 \(\dim E_\lambda\) 列的每一列都只是 \(\lambda\) 乘某个坐标向量；展开特征多项式时，因子 \((t-\lambda)^{\dim E_\lambda}\) 必然被提出来，代数重数就不会小于几何重数。

两者是否相等，决定了矩阵能不能凑够 \(n\) 个独立的特征向量。下一节的对角化正是以这个条件为前提。

\subsection{不同特征值的特征向量自动独立}

设 \(v_1,\ldots,v_k\) 分别对应互不相同的特征值 \(\lambda_1,\ldots,\lambda_k\)。这些向量一定是线性无关的。

两个向量的情形已经足够看出原理。若

\[
c_1v_1 + c_2v_2 = 0,
\]

两边左乘 \(A\)：

\[
c_1\lambda_1v_1 + c_2\lambda_2v_2 = 0.
\]

用第二式减去 \(\lambda_2\) 乘第一式：

\[
c_1(\lambda_1-\lambda_2)v_1 = 0.
\]

因为 \(v_1\neq0\) 且 \(\lambda_1\neq\lambda_2\)，必须 \(c_1=0\)，代回即得 \(c_2=0\)。一般情形可以归纳：每次用最后一项的特征值乘前一步的等式，相减消去最后一项，逐步推出所有系数为零。

推论很直接：如果一个 \(n\times n\) 矩阵有 \(n\) 个互不相同的特征值，它就有 \(n\) 个独立的特征向量，一定可以对角化。但不同特征值不是必要条件——单位矩阵只有一个特征值却满满地对角化。重复根本身不是问题，问题是重复根到底提供了多少个独立方向。

\subsection{迹和行列式怎样汇总特征值}

在复数范围内，把特征值按代数重数列为 \(\lambda_1,\ldots,\lambda_n\)，则

\[
\trace(A) = \lambda_1 + \cdots + \lambda_n,
\qquad
\det A = \lambda_1 \cdots \lambda_n.
\]

行列式汇总各特征方向的缩放乘积，迹汇总缩放之和。这两个关系可以从特征多项式的系数直接读出：\(\det(A-\lambda I)\) 的 \(\lambda^{n-1}\) 项系数含 \(-\trace(A)\)，常数项是 \(\det A\)。

它也解释了为什么 \(A\) 和 \(S^{-1}AS\) 有相同的特征值、迹和行列式：

\[
\det(S^{-1}AS-\lambda I)
= \det\bigl(S^{-1}(A-\lambda I)S\bigr)
= \det(A-\lambda I).
\]

因为 \(\det(S^{-1})\det(A-\lambda I)\det(S)=\det(A-\lambda I)\)。特征多项式不变，那些由它决定的量自然也不变。\(S^{-1}AS\) 形的矩阵与 \(A\) 相似——是同一个变换换了身基的描述，下一节会正式谈。

迹和行列式只是 \(n\) 个特征值的两个聚合统计量，知道它们不能恢复全部特征值。但有时这两个聚合量本身已经透露了重要结构。比如投影矩阵的特征值全部是 \(0\) 或 \(1\)——落在投影子空间中的向量完全不缩放（特征值 \(1\)），落在正交补中的向量被拍平（特征值 \(0\)）——所以投影矩阵的迹恰好等于它的秩，也等于投影子空间的维数。这个对应关系在后面处理统计学中的自由度和数值线性代数中的低秩近似时会反复出现。

\subsection{实矩阵也可能没有实特征方向}

平面旋转 \(90^\circ\) 的矩阵

\[
R = \begin{pmatrix}0&-1\\1&0\end{pmatrix}
\]

把每个非零向量都转了 \(90^\circ\)。在 \(\R^2\) 里，没有一条直线是旋转后仍落在自身上的。特征方程为

\[
\lambda^2+1=0,
\]

根为 \(\lambda=\pm i\)——没有实数根。

这不是计算出了故障。它忠实地反映了：实数平面上确实不存在保持方向不变的直线。把视野扩展到复数空间后，特征方向才出现——但这个解释本身要求引入复向量和复标量，后续的复矩阵单元会用共轭转置和酉基来处理这些情形。

延伸阅读：MIT 18.06：Eigenvalues and Eigenvectors。
